\section{Introduction}
\label{sec:intro}

Scene graph generation (SGG) represents an image with objects and their
directed relationships, providing a structured description for visual
reasoning \cite{xu2017scenegraph}. Open-vocabulary SGG (OVSGG) extends this
goal beyond the predicates observed during training. Generalizing relations is
particularly difficult: a predicate is determined jointly by the subject, the
object, and their interaction, while fine-grained relations often differ by
small visual cues.

Most OVSGG methods address this difficulty by improving the language side of
visual--text matching. Epic and RECODE construct cross-modal prompts or
composite descriptions \cite{yu2023epic,li2023recode}; RAHP, SDSGG, and APT
further incorporate entity, region, scene, or visual context into predicate
representations \cite{liu2025rahp,chen2024sdsgg,luo2026apt}. These methods
provide richer semantic targets, but a better target does not by itself ensure
a reliable transfer. Pointwise alignment constrains where each relation feature
should move, yet leaves its relative arrangement with other visual relations
unconstrained. It also offers no visual examples for predicates absent from
training.

We therefore view open-vocabulary relation recognition as a
geometry-preserving cross-modal distribution completion problem. A useful
transfer should satisfy two requirements: it should retain the visual
neighborhood structure that separates observed relation instances, and it
should provide visual support for unseen predicates rather than relying on
text-only extrapolation. This view shifts attention from designing ever richer
prompts to constraining the structure and visual coverage of the transfer.

Based on this view, we propose a relation-preserving modality-transfer
framework. A frozen CLIP triplet teacher \cite{radford2021clip} encodes each
subject--predicate--object combination and removes the dominant shared direction
from the resulting text bank. A learnable projector aligns visual relation
features with these triplet semantics, while visual structure preservation
(VSP) keeps pairwise similarities consistent before and after projection. We
further adapt synthesized prompting \cite{wang2023ship} to generate
relation-conditioned pseudo visual features for unseen predicates. At
inference, the resulting triplet-matching scores are fused with the original
predicate scores.

Experiments on Visual Genome and GQA validate these design choices. Our method
improves a strong recent baseline on all 30 reported Recall and
mean-Recall metrics, by 2.06 percentage points on average, and outperforms APT
on all 12 metrics under the matched Visual Genome protocol. Removing VSP lowers
every completed base/novel metric, whereas relation-conditioned synthesis
improves every novel metric. Our contributions are threefold: (1) we identify
structure preservation and missing novel visual support as two coupled problems
in open-vocabulary relation transfer; (2) we develop a unified framework that
preserves visual relation geometry while completing unseen visual support; and
(3) we demonstrate consistent improvements on two datasets and isolate the
roles of both aspects through ablation.
